\scanpage{12}\mspage{5}
Soit maintenant
\[
e_M\in(\det M)^*\tag*{(24)}
\]
une jauge de $M$, d'où un sous-ensemble
\[
S\widetilde F_M=\{\deleted{f}A\in\widetilde F_M\mid\alpha_A(\underline{\omega}_A)=\{\pm\deleted{\unsure{a|e}}\unsure{e}_M\}\}\tag*{(25)}
\]
Notons que $-\mathrm{id}_M$ applique $S\widetilde F_M$ dans lui-$=$, ce qui permet de définir
\[
SF_M=S\widetilde F_M/\{\pm\mathrm{id}_M\}\lhook\joinrel\longrightarrow\unsure{\mathfrak P|\mathscr P}_3(S_M)\tag*{(26)}
\]
On a encore, comme on voit tout de suite
\[
\widetilde S_M=\bigcup_{A\in S\widetilde F_M}A\tag*{(27)}
\]
Car si $x\in\widetilde S_M$ appartient à la base $(x,y)$, alors $\exists!\lambda\in\Lambda^*$, \deleted{\unsure{déterminé par}} \inserted{tel que} $x\wedge\lambda y=e_M$, donc $\{x,\lambda y,-(x+\lambda y)\}\in S\widetilde F_M$, d'où
\[
S_M=\bigcup_{f\in SF_M}f.\tag*{(28)}
\]
On pose
\[
\begin{aligned}
SA_M&=\{a\in\mathfrak P_2(S_M)\mid\exists f\in SF_M,\ a\subset f\}\\
&=\left\{a\in\mathfrak P_2(S_M)\ \middle|\
\begin{gathered}\exists(x,y)\in\mathrm{Rep}(M),\\ x\wedge y=e_M,\ a=\{\dot x,\dot y\}\end{gathered}\right\}\\
&\deleted{=\{a\in A_M}
\end{aligned}\tag*{(29)}
\]
Notons que
\[
SF_M\subset F_M,\quad SA_{\underline M}\subset A_M\tag*{(30)}
\]
\deleted{\illegible{1}} sont des sous-ensembles, qui peuvent \unsure{se|en}
% archive sheet number: 11
